\documentclass[11pt]{letter}

\usepackage[margin=0.5in]{geometry}
\usepackage{hyperref}

\begin{document}

\begin{letter}{Subsurface Dynamics}

\opening{Dear Hiring Team,}

I am writing to apply for the Data Integration Engineer (Real-Time Systems) position at Subsurface Dynamics. I recently completed my M.Sc. in Computer Science at the University of Calgary, and I am especially interested in this role because it sits close to the kind of work I enjoy most: making messy, high-volume operational data reliable enough that people can use it quickly and confidently.

My graduate research with BigGeo, Mitacs, and the University of Calgary focused on scalable data systems for Earth observation data. I built C++ and Python software for processing, storing, and retrieving large satellite datasets, automated SAR and optical imagery preprocessing, and developed DGGS-based encoding pipelines that normalized imagery into structured tensor files. This work required careful handling of imperfect source data, schema-like transformations, performance constraints, and debugging across data acquisition, processing, storage, and retrieval.

The ingestion and reliability focus of your platform is a strong match for that background. In my DGGS project, I designed data structures and retrieval workflows for multiresolution and time-series access, improved C++ processing performance from 233 seconds to 26 seconds for a 10-tensor workload, and reduced storage by 38 percent on scaled Canada-wide data. In my DEM super-resolution work, I built Python pipelines with Rasterio and Google Earth Engine to collect, preprocess, validate, and evaluate terrain data at scale. These projects taught me to care about correctness, repeatability, and operational clarity, not only about producing code that runs once.

I also bring backend engineering experience from my work as a Back-End Developer at Asr Gooyesh Pardaz, where I developed Django REST APIs, authentication features, PostgreSQL-backed service logic, and payment-service integrations for NLP, text-to-speech, and speech-to-text products. That experience maps well to your need for someone comfortable with APIs, JSON-style service data, database-backed records, and debugging issues across system boundaries.

What stands out to me about Subsurface Dynamics is the practical pressure of the problem: data is flowing during live frac operations, field decisions depend on it, and reliability matters immediately. I do not have direct SCADA or frac-domain experience yet, but I do have a strong foundation in Python, backend systems, data pipelines, time-series and geospatial data, Linux-based development, and careful debugging. I would be excited to learn the domain quickly and contribute to ingestion tooling, validation checks, mapping logic, and observability improvements that reduce manual intervention.

Thank you for considering my application. I am based in Calgary and would welcome the opportunity to discuss how my software, backend, and data-pipeline background could contribute to Subsurface Dynamics' real-time monitoring platform.

\closing{Sincerely,

Amir Mirzai Golpayegani

\href{mailto:amirgolpaa24@gmail.com}{amirgolpaa24@gmail.com}}

\end{letter}

\end{document}
